\paragraph{Scope of exact sampling}
The sampling rule is exact for the intended target distribution only when the supplied probabilities are correct. We restrict this claim to penalty-free sampling: presence and frequency penalties are zero, the repetition-penalty multiplier is one, min-$p$ is zero, and probability filters use matching tie conventions. The tree path ignores min-$p$, and its top-$k$ implementation retains exactly $k$ tokens, whereas native single-row decoding retains boundary ties. These differences must be absent or aligned even in the penalty-free regime. The numerical continuation result does not prove joint output-distribution equality.

The deployed and SWE-bench runs instead use presence penalty 1.0. Their logits processor constructs a flattened draft history, as for a chain, rather than each node's ancestry. A target row should include the path through its parent; a leaf continuation row should include the leaf itself. Flattening can penalize tokens from sibling branches and omit the current leaf token. The first token of a step has the correct penalty history; subsequent decisions can differ. We therefore describe these runs as departing from exact autoregressive sampling after the first token of each step, without asserting that every subsequent token changes.

\paragraph{Captured-row and sampler audit}
A passive capture of 40 tree steps on four requests contains 2,480 processed target and continuation rows, including padded slots. All match the flattened-chain history; only 63 coincide with the correct path history. Reconstructed probabilities from captured top-256 logits estimate node-level total variation above 0.01 on approximately 22\% of deeper decisions, with maximum 0.421. These estimates use rounded logits and depend on filter tie handling; they are not joint-sequence total variation.

The deployed target-only walk consumes target and continuation logits without a separate draft-probability tensor. Its CPU floating-point oracle is sampled 200,000 times on each of six captured steps, for 1.2 million walks. Tests compare emitted-token counts conditional on visiting a node with the supplied probability row, using an approximate chi-square $G$-test and a Bonferroni-adjusted per-step threshold of $10^{-3}$. Nodes require at least 2,000 visits; rare expected counts are pooled. No departure is detected in 54 node tests against the supplied rows. Comparison with corrected ancestry probabilities rejects five of six steps. All five planted fault types are detected in at least one step (Table~\ref{tab:sampling-faults}). Non-rejection supports this finite sampler check; it does not prove universal exactness or validate the deployed penalty histories.

\begin{table}[t]
\caption{Sampler-audit fault sensitivity on six captured steps. Each fault type is detected at least once; detection depends on which nodes the walk visits.}
\label{tab:sampling-faults}
\centering\small
\begin{tabular}{@{}lr@{}}
\toprule
Injected fault & Steps detected \\
\midrule
Temperature applied twice & 6/6 \\
Target rows shifted & 6/6 \\
Fixed residual draw & 5/6 \\
Unconditional acceptance & 3/6 \\
Correlated random draws & 1/6 \\
\bottomrule
\end{tabular}
\end{table}

An earlier ten-boot, 8,000-continuation output diagnostic failed its frozen temperature-control sensitivity requirement and remains uninformative. The tree sampler also seeds one process-wide generator with zero, so boots with identical request order reuse random sequences. They cannot be treated as independent sampling replicates. The probability-history audit identifies a specific deployed defect despite that uninformative output test.
